% SOURCE docs/release/2026-10-29/source/docs/writing-repair/2026-10-29/src/4ec6d032be31f60c.tex CHAR 0
\chapter{What a public experiment can tell a bank}
\label{gap29:public}

A useful prototype should take an intervention question through to a decision
that an economist can defend. Public data permits that exercise before a bank
commits staff time, customer records and money to a pilot. Its value depends on
which question it answers. A public experiment can test whether an analysis
preserves the population, intervention and uncertainty. It cannot establish
that the same intervention works for a bank's customers.

The National Supported Work demonstration provides a concrete example. In the
retained experimental subset, the estimated increase in 1978 earnings is about
\$1,794, with an approximate 95\% interval from \$479 to \$3,109. These are
the archive's earnings units, without conversion to today's purchasing power.
The result supports a positive average contrast under the stated design
assumptions. It does not establish that the gain exceeds an illustrative
\$1,000 cost. That distinction is exactly the distinction a banking proposal
must preserve when discussing a profitable offer or servicing intervention.

\section{The population is part of the question}
The author archive supplies separate files for 185 treated and 260 control
observations with earnings in 1974, 1975 and 1978. The original paper explains
that the subset is drawn from male experimental participants for whom 1974
earnings can be recovered. Selection uses pre-intervention information,
including assignment month and employment history. The authors explicitly
condition the retained experiment's interpretation on the original
randomisation being independent of those pre-intervention covariates.
\citep{gap29dehejia,gap29data}

The distinction matters. These 445 people are neither the full demonstration
nor a sample of bank customers. The public files do not contain the original
assignment logs or a complete site-by-month randomisation history. We retain
the paper's selection argument as a substantive premise. We do not infer it
from balance tests. Treatment assignment, rather than an independently
identified effect of treatment receipt, defines the contrast. Consistency,
absence of interference and preservation of the observed cohort are additional
premises.

For participant $i$, let $Y_i(1)$ and $Y_i(0)$ denote 1978 earnings under
assignment to the programme and control. The finite-population target is
\[
 \tau_S=\frac{1}{445}\sum_{i\in S}\{Y_i(1)-Y_i(0)\},
\]
where $S$ is the retained subset. The realised data reveal one potential
outcome per participant. The calculation therefore has uncertainty from treatment assignment
even though every row in the extract is used. Under the finite-population
randomisation model, this uncertainty concerns which potential outcome is
revealed for each participant. It neither observes individual
causal effects nor establishes an effect for an unobserved population.

The original data files contain all required variables without missing values.
Comparison with the previously retained teaching mirror finds a maximum
absolute difference of 0.00125 across matched fields, consistent with the
different numeric storage precision. We use the author's text files for the
reported calculations and retain their exact identities. The current archive
states attributable non-commercial terms; commercial redistribution or product
embedding requires a separate rights decision. The page was accessed after
the literature cutoff; its current bytes are not represented as a historical
snapshot. \citep{gap29data}

\section{Start with the simple contrast}
The primary estimator is the difference between treated and control means:
\[
 \widehat\tau_U=\overline Y_1-\overline Y_0,\qquad
 \widehat V_U=\frac{s_1^2}{185}+\frac{s_0^2}{260}.
\]
Here $s_a^2$ is the sample variance within assignment group $a$. Under the
complete-randomisation idealisation, the variance of the difference also
contains a subtraction for the unobserved variation in individual treatment
effects. Omitting that term provides the usual conservative variance
construction. The interval uses a normal approximation; it is not an exact
finite-sample statement. The actual site's assignment mechanism and selection
history remain qualifications on applying that idealisation.

Adjustment uses eight documented pre-intervention fields: age, education,
race and ethnicity indicators, marital status, degree status and earnings in
1974 and 1975. The comparison fixes this list before fitting; it does not
select variables to obtain a favourable result. For numerical conditioning,
each field is centred at its full-sample mean and divided by its full-sample
standard deviation.

The common-slope regression fits an intercept, assignment and these controls.
The interacted regression additionally includes every assignment-by-control
interaction. At the centred covariate origin, its assignment coefficient is
the difference between the two arm-specific fitted outcomes averaged over
the same pooled covariate distribution. In the design-based analysis of
\citet{gap29lin}, the fully interacted estimator has no worse asymptotic
precision than the unadjusted estimator under a fixed number of covariates,
bounded moments, a nonsingular limiting design and treatment proportions
bounded away from zero and one. This asymptotic result does not order the
estimated standard errors in one finite dataset.

For a design matrix $Z$, residuals $\widehat e_i$ and leverages $h_i$, the
reported regression covariance is
\[
 \widehat V_{HC2}=(Z'Z)^{-1}
 \left(\sum_i z_i z_i'\frac{\widehat e_i^2}{1-h_i}\right)(Z'Z)^{-1}.
\]
The HC3 sensitivity replaces $1-h_i$ by $(1-h_i)^2$. Neither adjustment
repairs a false causal model. Sandwich variance theory supplies asymptotic
support under the declared design conditions; small samples and influential
observations still warrant scrutiny. \citep{gap29lin}



% SOURCE docs/release/2026-10-29/source/docs/writing-repair/2026-10-29/src/0ab390ccab38d540.tex CHAR 0

\begin{table}[htbp]
\centering\small
\begin{tabular}{lrrl}
\toprule
Estimator & Estimate & Std. error & Approx. 95\% interval \\
\midrule
Difference in means & 1,794 & 671 & [479, 3,109] \\
Common-slope adjustment & 1,676 & 677 & [349, 3,003] \\
Fully interacted adjustment & 1,622 & 695 & [260, 2,983] \\
\bottomrule
\end{tabular}
\caption{New calculations on the 445 retained experimental observations.
Intervals use normal critical values and HC2 covariance, equivalent to
separate-arm variance for the unadjusted estimator. This single dataset
establishes neither finite-sample coverage nor estimator superiority.}
\end{table}


% SOURCE docs/release/2026-10-29/source/docs/writing-repair/2026-10-29/src/4ec6d032be31f60c.tex CHAR 5959


The adjusted point estimates are somewhat smaller, and their estimated
standard errors are somewhat larger. This dataset provides no basis for
declaring a more elaborate estimator superior. It also does not refute an
asymptotic theorem: estimated precision in a realised sample and repeated-
assignment variance are different objects. An honest software comparison
reports both the target and the uncertainty convention before discussing
numerical agreement.

\section{Precision, influence and the spending threshold}
About 30.8\% of observed 1978 earnings are zero; the largest is about
\$60,308. The largest leverage is 0.346 for the interacted fit. Its HC3
standard error is about 717, compared with 695 under HC2. Deleting one
observation at a time moves the unadjusted point estimate between about
1,501 and 1,929. This sensitivity is reported without selecting a preferred
sample or substituting a trimmed estimand.

The predeclared split by zero versus positive 1975 earnings gives similar
point estimates, about 1,711 and 1,691, but very different uncertainty. The
respective intervals are approximately [263, 3,159] and [-875, 4,258].
These are descriptive subgroup contrasts. Their comparison neither
establishes equal effects nor justifies choosing a subgroup after seeing
its significance. A targeted policy would require its own evaluation.

\begin{figure}[htbp]
\centering
\includegraphics[width=.94\textwidth]{scientific-gap-closure/2026-10-29/decision-intervals.pdf}
\caption{Estimated earnings contrasts and approximate 95\% intervals. The
dashed line is an illustrative cost of 1,000 in the archive's earnings units.
It is not an observed programme cost.}
\label{gap29:intervals}
\end{figure}

Let $c$ be a stipulated per-person cost in the same units. Subtracting $c$
from the estimated earnings contrast gives a deliberately limited accounting
comparison, $\widehat\tau-c$. At $c=1,000$, the unadjusted point comparison
is positive, about 794, but its interval is approximately [-521, 2,109].
A rule requiring the lower confidence bound to exceed cost therefore cannot
support spending at that threshold. At $c=2,000$, even the point comparison
is negative. Choosing the relevant rule requires a decision-maker's loss
function; statistical significance alone does not supply one.

Earnings gains are not automatically fiscal savings, bank revenue or social
surplus. Programme costs, displacement, persistence, discounting and the
distribution of benefits would all matter in a full economic appraisal.
The illustration asks whether the effect is large enough relative to an
explicit hurdle; it supplies no measured return on investment.

Sampling uncertainty is also distinct from a failure of the design premises.
If an analyst supplies a credible bound $B$ on the absolute discrepancy
between the estimated target and the target of interest, the interval can be
expanded from $[L,U]$ to $[L-B,U+B]$. With $B=500$, the unadjusted lower
endpoint falls to about -21. This is a conditional sensitivity calculation.
The data do not identify $B$, and the expansion is not a universal correction
for hidden confounding.

\section{What software agreement establishes}
The same rows and specifications were fitted independently with NumPy least
squares and R 4.1.2's \texttt{stats::lm}. R uses its QR-based fitting route;
the HC2 covariance was constructed explicitly rather than silently using a
default homoskedastic standard error. Point estimates, variances and standard
errors agree to within $10^{-6}$ in absolute units. The unadjusted covariance
also agrees with direct group-variance arithmetic, and the interacted
coefficient agrees with separate-arm standardisation. \citep{gap29r}

This result supports reusing existing linear-model machinery for this
specified calculation. It does not establish agreement with every causal
package, validate an R adapter, or qualify missingness, survey weights,
dependent records or other models. The proposal's integration should preserve
the choice of target and covariance explicitly. A language boundary is a
manageable engineering problem; an unrecorded change of statistical meaning
can invalidate the economic conclusion.

The fitted estimates and covariance summaries were then passed through the
existing library, command-line interface and local MCP service. Each returned
the same interval. In this construction the numerator is the estimated effect,
the denominator is the fixed value one, and the supplied covariance is
$\operatorname{diag}(n\widehat V,0)$. The service's Gaussian residual
inversion consequently reduces to the stated normal interval. This is an
executed path from public data to fitted summaries to consistent inference
and a local decision report. Raw-data estimation and report generation are
still local scripts, not newly implemented service endpoints.

Marking the assignment premise contradicted produces a refusal. Keeping it
unverified preserves conditional inference rather than certifying truth.
These behaviours make the worked example useful to an economist: uncertainty
and missing evidence survive the change of interface. They do not qualify a
bank platform or a sibling repository.

